% Imporditud: tools/import_eksperimendid.py
% Allikas: Eesti füüsikaolümpiaadi eksperimendi ülesannete
%   kogu (Rael Kalda, 2023), ülesanne Ü150, lahendus L150.
\setAuthor{Koit Timpmann}
\setRound{lõppvoor}
\setYear{2018}
\setNumber{P E2}
\setDifficulty{4}
\setTopic{elektriahelad}
\setVahendid{taskulambipirn alusel, lapik taskulambipatarei, 4 juhet, multimeeter (voolutugevuse mõõtmisel kasutage skaalat 0 A - 10 A)}
\setPunktid{12}
\setAllikas{Eksperimendi kogu Ü150, lahendus lk 111}
\setLahendusseis{kasitsi}

\prob{Hõõgniidi temperatuur}
Määrake töötava taskulambipirni hõõgniidi temperatuur. Toa temperatuur on 23 $^\circ$C ja volframi eritakistuse temperatuurikoefitsent $\alpha$ = 0,0044 1/K. Hõõgniidi soojuspaisumist ei ole tarvis arvestada.
\textit{Märkus:} Eritakistuse sõltuvust temperatuurist kirjeldab seos $\rho$ = $\rho_0$(1 + $\alpha$t).

\hint

\solu
\textit{Lahendus on kogumikus lk 111. Siia ei ole seda veel ümber kirjutatud, sest originaalis on tabel või joonis, mis PDF-i tekstikihist loetavalt välja ei tule.}
\probend
